\section{GlyphChain and Legal production-release evidence gate}
\label{sec:glyphchain-legal-release-evidence-gate}

The production decision is now protected by a fail-closed evidence evaluator
covering every unfinished GlyphChain and Legal release task.  The evaluator does
not manufacture evidence, appoint reviewers or approve a release merely because
the software is installed.  It verifies that independently configured
authorities signed production evidence for the exact release under review and
keeps the release blocked until all thirty-four gates are satisfied.

\subsection{Canonical task register}

The register contains exactly thirty-four identifiers divided into the six
remaining workstreams:

\begin{center}
\begin{tabular}{lll}
\hline
\textbf{Prefix} & \textbf{Workstream} & \textbf{Task count} \\
\hline
\texttt{LD} & Live signing delivery & 5 \\
\texttt{TR} & Independent quorum assurance publication & 3 \\
\texttt{NET} & Distributed-network production validation & 5 \\
\texttt{SEC} & Production security and recovery & 4 \\
\texttt{LEGAL} & Qualified Legal content & 7 \\
\texttt{REL} & Controlled deployment and release & 10 \\
\hline
& \textbf{Total} & \textbf{34} \\
\hline
\end{tabular}
\end{center}

Each identifier has a fixed description, required authority roles and an
independence requirement where appropriate.  Unknown identifiers do not extend
the register.  Duplicate entries are reported and block approval.

\subsection{Exact-release binding}

Every evidence body contains the SHA-256 digest of the exact release manifest.
For release bytes $M$, the binding is

\begin{equation}
  h_M = \operatorname{SHA256}(M).
\end{equation}

An entry for another build, branch, package or earlier candidate cannot satisfy
a task, even when its underlying exercise passed.  Human-readable release names
are not accepted as substitutes for $h_M$.

The signed evidence body includes:

\begin{itemize}
  \item schema and task identifier;
  \item exact release-manifest digest;
  \item the literal \texttt{production} environment;
  \item a \texttt{PASS} result;
  \item observation time and retained statement;
  \item one or more uniquely identified, task-specific artifact SHA-256 digests;
        and
  \item the attestations required by the task's authority policy.
\end{itemize}

Each task declares the artifact classes it requires.  An unrelated report with a
valid digest cannot satisfy the gate; all required classes must be present and
artifact identifiers must be unique inside the evidence body.  Local, test or
staging evidence is rejected by the production gate.  A task
also remains open when its result is incomplete, its retained artifact list is
empty or malformed, or its statement and observation time are absent.

\subsection{Role-scoped Ed25519 attestations}

An authority signs a domain-separated canonical body:

\begin{equation}
  \sigma_i = \operatorname{Ed25519.Sign}_{sk_i}\!\left(
    \operatorname{canon}(D,E)
  \right),
\end{equation}

where $D$ is
\texttt{TESSARIS\_GLYPHCHAIN\_LEGAL\_RELEASE\_EVIDENCE\_V1} and $E$ is the
complete evidence body.  The configured trust-root inventory supplies the
matching public key, permitted roles and whether the authority is independent
of the development team.  A non-empty inventory is accepted only when its
governance signature verifies against a separately supplied Ed25519 authority
key and its exact file digest matches the independently supplied current pin.

The role set distinguishes delivery operators, security reviewers, independent
trust publishers and verifiers, network and security operators, application
and consensus security reviewers, qualified Legal reviewers, Legal product
owners, deployment operators, pilot owners, metrics reviewers, release
publishers and final release approvers.

A valid signature from the wrong role is insufficient.  For gates requiring
multiple roles, each role must be represented by a distinct signer.  One person
or key holding several labels cannot satisfy separation of duties for that
task.  When independence is required, at least one of the signers actually
satisfying a required role must be marked independent in the separately
controlled trust-root inventory.

\subsection{Release-decision semantics}

For task $t$, define $G_t=1$ only when its body, release binding, artifacts,
required roles, signatures, independence and separation-of-duties checks all
pass.  Production authorization is computed as

\begin{equation}
  G_{\mathrm{release}} =
  \left(\bigwedge_{t \in T} G_t\right)
  \land (U=\varnothing)
  \land (D=\varnothing),
\end{equation}

where $T$ is the fixed thirty-four-task register, $U$ is the set of unknown task
identifiers and $D$ is the set of duplicates.  There is no percentage threshold,
majority override or silent waiver.

If the conjunction is false, the result is
\texttt{BLOCKED\_EXTERNAL\_GATES} and
\texttt{production\_release\_authorized=false}.  Only a complete evidence set
can yield \texttt{ALL\_RELEASE\_GATES\_EVIDENCED}.  The evaluator verifies the
configured evidence; it does not independently determine whether a person is a
qualified lawyer, reproduce an external provider event or replace the signed
human approval.

\subsection{Evidence inventory and trust administration}

The repository includes deliberately empty inventory templates.  The evidence
inventory contains no self-declared passing records, and the trust-root
inventory authorises no keys.  Production operators must populate them only
through the approved external governance process.  Private signing keys must
never be stored in either file or in the tracked source tree.

Trust roots and evidence are deployment records rather than source-code proof.
They are mutable only through controlled operations and are intentionally not
folded into the source manifest that they are intended to approve.  This avoids
a circular construction in which changing the approval evidence changes the
release digest it signs.

\subsection{Integrity-checksummed readiness report}

The evaluator produces an owner-only, exclusive-create JSON report.  It records
the release digest, passed and unmet identifiers, unknowns, duplicates,
verified role-to-signer mappings, artifact counts and per-task errors.  It does
not store private keys or private Legal content.

The report also binds canonical SHA-256 digests of the complete evidence-entry
sequence and trust-root map.  Each task records the exact evidence-body digest,
artifact-manifest digest and attesting signer identifiers.  This makes the
decision reproducible against the signed source evidence and exposes inventory
substitution.  The report's own hash is an integrity checksum, not an independent
signature; authenticity derives from re-evaluating the retained Ed25519 evidence
against the separately administered trust roots.

For report body $R$, integrity is bound by

\begin{equation}
  h_R = \operatorname{SHA256}\!\left(
    \operatorname{canon}(R \setminus \{\texttt{report\_sha256}\})
  \right).
\end{equation}

The report writer refuses symbolic-link destinations and existing files.  A
change to any task result, signer mapping, error, release digest or claim
boundary invalidates $h_R$.

\subsection{Implemented components and verification}

The release policy and verifier are implemented in
\texttt{backend/modules/aion\_business/runtime/glyphchain\_release\_readiness.py}.
The command-line evaluator is
\texttt{backend/scripts/evaluate\_glyphchain\_legal\_release\_readiness.py}.
The regression set is
\texttt{backend/tests/workflow\_capsules/test\_aion\_legal\_glyphchain\_release\_readiness.py}.

Direct regressions prove: a blank inventory leaves all thirty-four tasks
open; valid evidence closes only its exact task; wrong release digests,
signatures and roles fail closed; independent gates reject non-independent
authority, including a textual \texttt{"false"} flag; multi-role gates use a
complete distinct-signer assignment rather than a greedy role selection;
malformed, timezone-free and future observation times fail closed; a complete
correctly signed synthetic inventory can satisfy the conjunction; exact evidence
and trust inventories are digest-bound; and unknown, duplicate, malformed or
altered evidence blocks the release.  The command-line evaluator also validates
both inventory schemas and imports without initialising unrelated application
runtimes.

The current retained readiness run uses the empty external-evidence and
trust-root inventories.  It correctly reports zero passed tasks, thirty-four
unmet tasks, \texttt{BLOCKED\_EXTERNAL\_GATES} and
\texttt{production\_release\_authorized=false}.  This is the accurate release
state until live providers, independent publication, representative networks,
security reviews, qualified Legal review and the controlled pilot produce real
signed evidence.
